\documentclass[12pt,a4paper]{article}
\usepackage[a4paper,margin=18mm]{geometry}
\usepackage[T2A]{fontenc}
\usepackage[utf8]{inputenc}
\usepackage[russian]{babel}
\usepackage{amsmath,amssymb,mathtools}
\usepackage{array,booktabs,longtable}
\usepackage{xcolor}
\usepackage{hyperref}
\usepackage{tikz}

\hypersetup{
  colorlinks=true,
  linkcolor=blue!55!black,
  urlcolor=blue!55!black
}
\setlength{\parindent}{0pt}
\setlength{\parskip}{0.55em}
\setlength{\tabcolsep}{3pt}
\renewcommand{\arraystretch}{1.18}

\newcommand{\errlabel}[1]{\textcolor{red}{Ошибка:} #1}
\newcommand{\keyidea}[1]{\textcolor{blue!60!black}{Ключевая идея:} #1}
\newcommand{\resultlabel}[1]{\textcolor{teal!60!black}{Итог:} #1}

\begin{document}

\begin{titlepage}
\centering
\vspace*{0.22\textheight}
{\LARGE\bfseries Ревью домашней работы\par}
\vspace{2.5em}
{\large Автор: Лёвин Артём Александрович\par}
\vspace{1em}
{\large 7 сентября 2026 года\par}
\vfill
\begin{tikzpicture}[x=1cm,y=1cm]
  \draw[blue!35!black, line width=0.7pt]
    (0,0) .. controls (2,0.45) and (4,-0.45) .. (6,0)
          .. controls (8,0.45) and (10,-0.45) .. (12,0);
\end{tikzpicture}
\end{titlepage}

\newpage
\section*{Сводная таблица}
\small
\begin{longtable}{>{\centering\arraybackslash}p{0.07\linewidth}p{0.18\linewidth}p{0.28\linewidth}p{0.19\linewidth}p{0.19\linewidth}}
\toprule
\textbf{№} & \textbf{Тема} & \textbf{Суть ошибки} & \textbf{Ваш ответ} & \textbf{Правильный ответ}\\
\midrule
\endfirsthead
\toprule
\textbf{№} & \textbf{Тема} & \textbf{Суть ошибки} & \textbf{Ваш ответ} & \textbf{Правильный ответ}\\
\midrule
\endhead
\hyperref[task:4]{4} & Касательная, диаметр, дуги & Для данных задания №4 отдельное решение не представлено: разобрана конфигурация с дугой $132^\circ$, относящаяся к №3. & не указан & $\angle ACO=14^\circ$ двумя способами\\
\midrule
\hyperref[task:5]{5} & Равные хорды и дуги & Из равенства сумм длин хорд сделан вывод о равенстве дуг. Такой переход неверен. & $AC=CE$ & $AC=CE$\\
\midrule
\hyperref[task:6]{6} & Вписанный четырёхугольник, трапеция, Птолемей & Есть чертёж, но нет доказательства $AD\parallel BC$ и вычисления диагоналей. & не указан & $AD\parallel BC$, $AC=BD=11$\\
\midrule
\hyperref[task:7]{7} & Равнобедренные трапеции в окружности, Птолемей & Равнобедренная трапеция заявлена без корректного обоснования; после $AC=BD=\sqrt{44}$ решение не доведено до $BE$. & не указан & $BE=\dfrac{19}{3}$\\
\midrule
\hyperref[task:10]{10} & Равные хорды и дуги, Птолемей & Повторён неверный переход от суммы длин хорд к дугам; затем получен необоснованный итог $BE=8$. & $BE=8$ & $BE=7$\\
\bottomrule
\end{longtable}
\normalsize

\newpage
\vspace*{\fill}
\begin{center}
{\LARGE\bfseries Разбор ошибок}
\end{center}
\vspace*{\fill}

\newpage
\phantomsection\label{task:4}
\section*{Задание 4}

\textbf{Текст задания.} В той же конфигурации: из точки $C$ к окружности проведена касательная $CA$, прямая $CO$ проходит через центр окружности и пересекает окружность в точках $B$ и $D$, точки расположены в порядке $C,B,O,D$. Дуга $AD$ равна $104^\circ$. Найти $\angle ACO$ двумя способами: через прямоугольный треугольник и через формулу угла между касательной и секущей.

\textbf{Ваше решение.} Отдельного решения для дуги $104^\circ$ не представлено. На странице перед заданием №5 выполнены два способа вычисления для дуги $132^\circ$ и получено $42^\circ$; это соответствует данным задания №3, а не №4.

\textbf{Ваш ответ.} не указан.

\errlabel{задание №4 не решено по своим данным.}

\textbf{Где именно возникает ошибка.} Для задания №4 нет последнего корректного вычислительного шага: работа с числом $104^\circ$ не начинается. Первый неполный шаг --- отсутствие перехода от условия $\widehat{AD}=104^\circ$ к дуге $AB$ или к центральному углу $\angle AOB$.

\textbf{Почему это неверно.} В задачах одной конфигурации нельзя переносить численный ответ с предыдущего пункта. Здесь вся величина искомого угла зависит от меры дуги $AD$. При $132^\circ$ получается одно значение, при $104^\circ$ --- другое. Поэтому даже правильный метод, применённый к числу из соседнего задания, не является решением задания №4.

\keyidea{точки $B$ и $D$ --- концы диаметра, поэтому дуги $BA$ и $AD$ на одной полуокружности в сумме дают $180^\circ$. Кроме того, радиус $OA$ перпендикулярен касательной $AC$.}

\textbf{Исправление от первого необходимого шага.} Сначала используем именно данное число:
\[
\widehat{AB}=180^\circ-\widehat{AD}=180^\circ-104^\circ=76^\circ.
\]
Значит, центральный угол
\[
\angle AOB=76^\circ.
\]
Так как точки $C,B,O$ лежат на одной прямой и лучи $OB$ и $OC$ совпадают, то
\[
\angle AOC=76^\circ.
\]

\textbf{Полное правильное решение, способ 1: через прямоугольный треугольник.}

Радиус, проведённый в точку касания, перпендикулярен касательной:
\[
OA\perp AC.
\]
Следовательно, треугольник $AOC$ прямоугольный при $A$, и его острые углы в сумме дают $90^\circ$:
\[
\angle ACO=90^\circ-\angle AOC=90^\circ-76^\circ=14^\circ.
\]

\textbf{Полное правильное решение, способ 2: угол между касательной и секущей.}

Угол между касательной $CA$ и секущей $CBD$, проведёнными из одной внешней точки $C$, равен половине разности соответствующих дуг:
\[
\angle ACO=\frac12\left(\widehat{AD}-\widehat{AB}\right).
\]
Подставляем найденную дугу $AB$:
\[
\angle ACO=\frac12(104^\circ-76^\circ)=\frac12\cdot 28^\circ=14^\circ.
\]

\textbf{Проверка.} Оба независимых школьных способа дали один и тот же результат $14^\circ$. Это согласуется с конфигурацией: угол при внешней точке $C$ должен быть острым.

\resultlabel{$\angle ACO=14^\circ$.}

\textbf{Задача на закрепление.} В той же конфигурации точки расположены в порядке $C,B,O,D$, $CA$ --- касательная, $BD$ --- диаметр, а дуга $AD$ равна $118^\circ$. Найдите $\angle ACO$ двумя способами: через прямоугольный треугольник $AOC$ и через формулу угла между касательной и секущей.

\newpage
\phantomsection\label{task:5}
\section*{Задание 5}

\textbf{Текст задания.} Пятиугольник $ABCDE$ вписан в окружность. Известно:
\[
AB=CD,\qquad BC=DE.
\]
Докажите, что
\[
AC=CE.
\]

\textbf{Ваше решение.} В работе использована цепочка вида
\[
AB+BC=CD+DE=13,
\]
после чего сделан вывод, что дуги $ABC$ и $CDE$ равны, а затем --- что $AC=CE$. При этом в исходном условии задания №5 чисел $5$, $8$ и $12$ нет.

\textbf{Ваш ответ.} $AC=CE$.

\errlabel{равенство сумм длин хорд ошибочно используется как основание для равенства дуг.}

\textbf{Где именно возникает ошибка.} Последний корректный по смыслу факт --- из $AB=CD$ и $BC=DE$ действительно следует
\[
AB+BC=CD+DE
\]
как равенство \emph{длин отрезков}. Первый неверный переход --- считать, что из этого автоматически следует
\[
\widehat{ABC}=\widehat{CDE}.
\]

\textbf{Почему этот шаг неверен.} Длина хорды и градусная мера стягиваемой ею дуги связаны нелинейно. Поэтому сумма длин двух хорд не является мерой суммы двух дуг. Например, даже если $AB+BC=CD+DE$, это само по себе не позволяет сравнивать дуги $AB+BC$ и $CD+DE$. Нужное здесь свойство другое: \emph{равные хорды одной окружности стягивают равные дуги}. Его надо применять к каждой паре равных хорд отдельно.

Кроме того, числовая запись $=13$ относится к изменённым данным и не требуется для доказательства исходного утверждения. Задание №5 является общим: доказательство должно работать без подстановки конкретных длин.

\keyidea{сначала перейти от каждой пары равных хорд к равным дугам, а уже затем сложить равные дуги.}

\textbf{Исправление от последнего правильного шага.} Из
\[
AB=CD
\]
получаем
\[
\widehat{AB}=\widehat{CD}.
\]
Из
\[
BC=DE
\]
получаем
\[
\widehat{BC}=\widehat{DE}.
\]
Теперь складываем \emph{меры дуг}:
\[
\widehat{ABC}=\widehat{AB}+\widehat{BC}
=\widehat{CD}+\widehat{DE}
=\widehat{CDE}.
\]

\textbf{Полное правильное решение.} Так как равные хорды одной окружности стягивают равные дуги,
\[
AB=CD\Rightarrow \widehat{AB}=\widehat{CD},
\]
\[
BC=DE\Rightarrow \widehat{BC}=\widehat{DE}.
\]
Следовательно,
\[
\widehat{ABC}=\widehat{AB}+\widehat{BC}
=\widehat{CD}+\widehat{DE}
=\widehat{CDE}.
\]
Хорда $AC$ стягивает дугу $ABC$, а хорда $CE$ стягивает дугу $CDE$. Равным дугам одной окружности соответствуют равные хорды, поэтому
\[
AC=CE.
\]

\textbf{Проверка.} В доказательстве использованы только два данных равенства $AB=CD$ и $BC=DE$ и стандартное свойство равных хорд и дуг. Никакие дополнительные числовые условия не нужны.

\resultlabel{доказано, что $AC=CE$.}

\textbf{Задача на закрепление.} Пятиугольник $PQRST$ вписан в окружность. Известно, что $PQ=RS$ и $QR=ST$. Докажите, что $PR=RT$.

\newpage
\phantomsection\label{task:6}
\section*{Задание 6}

\textbf{Текст задания.} Четырёхугольник $ABCD$ вписан в окружность, причём
\[
AB=CD=5,\qquad BC=8,\qquad AD=12.
\]
Докажите, что
\[
AD\parallel BC,
\]
а затем найдите $AC$ и $BD$.

\textbf{Ваше решение.} Для задания выполнен чертёж с отмеченными длинами, однако доказательство параллельности и вычисления диагоналей не записаны.

\textbf{Ваш ответ.} не указан.

\errlabel{решение не доведено дальше чертежа; обязательные доказательство и вычисления отсутствуют.}

\textbf{Где именно возникает ошибка.} Последний корректный элемент --- перенос данных $AB=CD=5$, $BC=8$, $AD=12$ на чертёж. Первый неполный шаг --- отсутствие доказательства $AD\parallel BC$. Без него нельзя обоснованно объявлять четырёхугольник равнобедренной трапецией и использовать равенство диагоналей.

\textbf{Почему это важно.} Равенство двух сторон $AB=CD$ само по себе не делает произвольный четырёхугольник трапецией. Здесь решающую роль играет то, что все четыре вершины лежат на одной окружности: равные хорды $AB$ и $CD$ стягивают равные дуги, а равные дуги дают равные вписанные углы. Именно из этих углов и получается параллельность оснований.

\keyidea{доказать $AD\parallel BC$ через равные вписанные углы $\angle BCA$ и $\angle CAD$, после чего использовать свойство диагоналей равнобедренной трапеции и теорему Птолемея.}

\textbf{Исправление от последнего правильного шага.} Так как
\[
AB=CD,
\]
то равны соответствующие дуги:
\[
\widehat{AB}=\widehat{CD}.
\]
Вписанный угол $\angle BCA$ опирается на дугу $BA$, а $\angle CAD$ --- на дугу $CD$, поэтому
\[
\angle BCA=\angle CAD.
\]
Эти углы являются накрест лежащими при секущей $AC$, значит
\[
AD\parallel BC.
\]

\textbf{Полное правильное решение.}

Мы уже доказали $AD\parallel BC$. Поскольку боковые стороны
\[
AB=CD=5,
\]
четырёхугольник $ABCD$ является равнобедренной трапецией. У равнобедренной трапеции диагонали равны:
\[
AC=BD.
\]
Обозначим
\[
AC=BD=x.
\]
Так как четырёхугольник вписан в окружность, применяем теорему Птолемея:
\[
AC\cdot BD=AB\cdot CD+BC\cdot AD.
\]
Подставляем данные:
\[
x^2=5\cdot5+8\cdot12=25+96=121.
\]
Длина положительна, поэтому
\[
x=11.
\]
Следовательно,
\[
AC=BD=11.
\]

\textbf{Проверка.} Теорема Птолемея даёт
\[
11\cdot11=121,
\]
а правая часть равна
\[
5\cdot5+8\cdot12=25+96=121.
\]
Равенство выполняется.

\resultlabel{$AD\parallel BC$, $AC=BD=11$.}

\textbf{Задача на закрепление.} Четырёхугольник $KLMN$ вписан в окружность. Известно, что $KL=MN=3$, $LM=5$, $KN=8$. Докажите, что $KN\parallel LM$, а затем найдите диагонали $KM$ и $LN$.

\newpage
\phantomsection\label{task:7}
\section*{Задание 7}

\textbf{Текст задания.} Пятиугольник $ABCDE$ вписан в окружность. Дано:
\[
AB=CD=3,\qquad BC=DE=5,\qquad AD=7.
\]
Найдите $BE$.

\textbf{Ваше решение.} В работе записано, что $ABCD$ --- равнобедренная трапеция, затем применена теорема Птолемея:
\[
5\cdot7+3\cdot3=y^2,
\]
\[
35+9=y^2,
\]
\[
y=\sqrt{44}.
\]
Также отмечено равенство диагоналей $AC=BD$. До вычисления $BE$ решение не доведено.

\textbf{Ваш ответ.} не указан.

\errlabel{параллельность оснований трапеции не доказана, а после верного промежуточного вычисления решение остановлено до искомой величины $BE$.}

\textbf{Где именно возникает ошибка.} Первый логический пробел появляется в утверждении, что $ABCD$ является равнобедренной трапецией только из-за $AB=CD$. Теорема Птолемея не доказывает параллельность сторон. Однако в данной вписанной конфигурации нужная трапеция действительно получается, если сначала доказать параллельность через равные дуги и вписанные углы.

После этого вычисление
\[
AC=BD=\sqrt{44}=2\sqrt{11}
\]
корректно. Последний корректный вычислительный шаг именно этот. Первый неполный шаг после него --- отсутствие перехода к вписанному четырёхугольнику $BCDE$ и вычислению $BE$.

\textbf{Почему это неверно или неполно.} Искомая величина --- $BE$, поэтому нахождение $AC$ и $BD$ является только промежуточным этапом. Нужно использовать вторую пару равных хорд $BC=DE$, чтобы получить ещё одну равнобедренную трапецию $BCDE$, а затем ещё раз применить теорему Птолемея.

\keyidea{в одном вписанном пятиугольнике последовательно использовать две равнобедренные трапеции: $ABCD$ из $AB=CD$ и $BCDE$ из $BC=DE$.}

\textbf{Исправление от первого логического пробела.}

Из $AB=CD$ следуют равные дуги $AB$ и $CD$. Поэтому
\[
\angle BCA=\angle CAD,
\]
и, следовательно,
\[
AD\parallel BC.
\]
Значит, $ABCD$ --- равнобедренная трапеция и
\[
AC=BD.
\]
По теореме Птолемея:
\[
AC\cdot BD=AB\cdot CD+BC\cdot AD.
\]
Отсюда
\[
AC^2=3\cdot3+5\cdot7=44,
\]
то есть
\[
AC=BD=2\sqrt{11}.
\]

Теперь рассматриваем вписанный четырёхугольник $BCDE$. Так как $BC=DE$, равны дуги $BC$ и $DE$. Вписанные углы
\[
\angle EBD \quad\text{и}\quad \angle BDC
\]
опираются соответственно на эти равные дуги, значит они равны. Это накрест лежащие углы при секущей $BD$, поэтому
\[
BE\parallel CD.
\]
Следовательно, $BCDE$ --- равнобедренная трапеция, а её диагонали равны:
\[
BD=CE=2\sqrt{11}.
\]

\textbf{Полное правильное решение.} Применяем теорему Птолемея к вписанному четырёхугольнику $BCDE$:
\[
BD\cdot CE=BC\cdot DE+BE\cdot CD.
\]
Подставляем известные значения:
\[
(2\sqrt{11})(2\sqrt{11})=5\cdot5+3\cdot BE.
\]
Левая часть равна $44$:
\[
44=25+3BE.
\]
Тогда
\[
3BE=19,
\]
\[
BE=\frac{19}{3}.
\]

\textbf{Проверка.} Подставим найденное значение в теорему Птолемея для $BCDE$:
\[
5\cdot5+3\cdot\frac{19}{3}=25+19=44,
\]
а
\[
BD\cdot CE=(2\sqrt{11})^2=44.
\]
Обе части совпадают.

\resultlabel{$BE=\dfrac{19}{3}$.}

\textbf{Задача на закрепление.} Пятиугольник $PQRST$ вписан в окружность. Известно, что $PQ=RS=4$, $QR=ST=5$, $PS=9$. Найдите $QT$.

\newpage
\phantomsection\label{task:10}
\section*{Задание 10}

\textbf{Текст задания.} Пятиугольник $ABCDE$ вписан в окружность. Известно:
\[
AB=CD=4,\qquad BC=DE=6,\qquad AD=8.
\]
1) Докажите, что $AC=CE$. 2) Найдите $BE$.

\textbf{Ваше решение.} В первой части использовано равенство
\[
AB+BC=10,\qquad ED+DC=10,
\]
после чего сделан вывод о равенстве дуг $ABC$ и $EDC$ и, следовательно, $AC=EC$. Во второй части для $ABCD$ применена теорема Птолемея:
\[
6\cdot8+4\cdot4=y^2,
\]
откуда
\[
y=8,
\]
и записано
\[
BD=AC=EC=8.
\]
Затем без корректного вычисления по четырёхугольнику $BCDE$ получен итог $BE=8$.

\textbf{Ваш ответ.} $BE=8$.

\errlabel{сначала сумма длин хорд ошибочно отождествляется с мерой дуги; затем пропущено необходимое применение теоремы Птолемея к $BCDE$, из-за чего итог $BE$ неверен.}

\textbf{Где именно возникает ошибка.} Первый неверный шаг находится уже в доказательстве пункта 1: из
\[
AB+BC=ED+DC
\]
как равенства длин отрезков нельзя делать вывод
\[
\widehat{ABC}=\widehat{EDC}.
\]
Правильный вывод должен строиться по отдельным равенствам хорд $AB=CD$ и $BC=DE$.

После правильного вычисления
\[
AC=BD=8
\]
вторая ошибка --- переход к ответу $BE=8$ без уравнения, связывающего $BE$ с известными сторонами и диагоналями четырёхугольника $BCDE$.

\textbf{Почему это неверно.} Равные \emph{суммы длин хорд} не определяют равные суммы дуг. Но равные \emph{отдельные хорды} одной окружности действительно стягивают равные дуги. Во второй части равенство $BD=CE=8$ ещё не означает $BE=8$: $BE$ --- основание другой равнобедренной трапеции, а не её диагональ. Для его нахождения нужно применить теорему Птолемея к $BCDE$.

\keyidea{пункт 1 доказать через равные дуги по отдельным парам хорд; пункт 2 решить двумя последовательными применениями структуры равнобедренной трапеции и теоремы Птолемея.}

\textbf{Исправление первой ошибки.} Из $AB=CD$ следует
\[
\widehat{AB}=\widehat{CD},
\]
а из $BC=DE$ следует
\[
\widehat{BC}=\widehat{DE}.
\]
Поэтому
\[
\widehat{ABC}=\widehat{AB}+\widehat{BC}
=\widehat{CD}+\widehat{DE}
=\widehat{CDE}.
\]
Следовательно, стягивающие эти дуги хорды равны:
\[
AC=CE.
\]

\textbf{Полное правильное решение.}

Сначала найдём $AC$ и $BD$. Из $AB=CD$ тем же способом, что в задании №6, получаем
\[
AD\parallel BC.
\]
Следовательно, $ABCD$ --- равнобедренная трапеция и
\[
AC=BD.
\]
Пусть
\[
AC=BD=x.
\]
По теореме Птолемея для $ABCD$:
\[
x^2=AB\cdot CD+BC\cdot AD
=4\cdot4+6\cdot8
=16+48=64.
\]
Значит,
\[
x=8,
\]
то есть
\[
AC=BD=8.
\]
Из доказанного в пункте 1 имеем
\[
CE=AC=8.
\]

Теперь применяем теорему Птолемея к вписанному четырёхугольнику $BCDE$:
\[
BD\cdot CE=BC\cdot DE+CD\cdot BE.
\]
Подставляем числа:
\[
8\cdot8=6\cdot6+4\cdot BE.
\]
Получаем
\[
64=36+4BE,
\]
\[
4BE=28,
\]
\[
BE=7.
\]

\textbf{Проверка.} Для $BCDE$:
\[
BC\cdot DE+CD\cdot BE=6\cdot6+4\cdot7=36+28=64,
\]
а
\[
BD\cdot CE=8\cdot8=64.
\]
Теорема Птолемея выполняется.

\resultlabel{доказано $AC=CE$; $BE=7$.}

\textbf{Задача на закрепление.} Пятиугольник $KLMNP$ вписан в окружность. Известно, что $KL=MN=5$, $LM=NP=7$, $KN=10$. 1) Докажите, что $KM=MP$. 2) Найдите $LP$.

\end{document}
